\begin{abstract}
\textbf{Background}: Identifying prodromal individuals at high risk for phenoconversion to clinical Parkinson's disease (PD) enables early intervention and preventive trial enrollment. We developed and validated deep neural survival models for personalized conversion risk prediction.

\textbf{Methods}: We analyzed 412 prodromal participants (REM sleep behavior disorder and/or hyposmia, no motor PD) from the Parkinson's Progression Markers Initiative, followed 5.1$\pm$1.6 years. Phenoconversion was defined as UPDRS-III$\geq$15 or clinical PD diagnosis. We compared Cox proportional hazards with Deep Neural Survival (\deepsurv) models using baseline clinical features, biomarkers (DaTscan, CSF $\alpha$-synuclein), and genetic variants. Models were validated via time-dependent AUC and calibration.

\textbf{Findings}: 75/412 participants (18.2\%) phenoconverted (median time-to-conversion 3.8 years). \deepsurv{} outperformed Cox models (C-index 0.76 vs. 0.68, $p=0.003$; 5-year AUC 0.78 vs. 0.71). Baseline UPDRS-III dominated risk prediction (hazard ratio 3.8 per 5-point increase, 95\% CI 2.3-6.2), followed by RBD (HR 3.2, 1.8-5.7), hyposmia (HR 2.4, 1.4-4.1), and GBA mutations (HR 3.5, 1.9-6.4). High-risk individuals (top tertile predicted risk, 31\% 5-year conversion) vs. low-risk (4\% conversion) achieved 89\% sensitivity, 78\% specificity. Enriching prodromal trials with high-risk individuals reduced required sample size by 42\% for neuroprotective interventions.

\textbf{Interpretation}: Deep neural survival modeling enables personalized prodromal-to-clinical PD conversion risk prediction from baseline assessments, with implications for early intervention trials and clinical counseling. The high-risk stratum identifies an enrichable population for disease modification studies.

\textbf{Funding}: Michael J. Fox Foundation, NINDS.
\end{abstract}
